\section{Introduction}
\label{sec:1}

Une perturbation logistique, qu'il s'agisse de la panne d'un fournisseur, de la fermeture d'un entrepôt ou d'un pic de demande, produit souvent ses effets pendant plusieurs semaines. Ce qu'elle coûte dépend de la sévérité du choc, et tout autant de la rapidité et de la pertinence avec lesquelles les gestionnaires réagissent \citep{sheffi2005book,ivanov2019ijpr}. Un réseau robuste absorbe le choc sans se dégrader ; lorsque ce n'est pas possible, sa résilience, c'est-à-dire sa capacité à limiter la perte puis à se rétablir, se joue en grande partie dans les décisions prises pendant la crise.

Les stratégies d'atténuation et de contingence sont bien répertoriées \citep{tomlin2006,chopra2004}, et la résilience est classiquement représentée par une courbe de perte et de rétablissement \citep{bruneau2003}. Il reste difficile, en revanche, de répondre à deux questions qu'un gestionnaire se pose en pratique. Face à une perturbation donnée, quelle réaction a le plus de chances d'atteindre le niveau de résilience visé ? Après un rétablissement médiocre, la perte tenait-elle au choc, au réseau ou à la réaction ? Pour y répondre, il faut des observations qui relient une perturbation, les décisions prises et l'évolution du taux de service qui s'ensuit. Les données d'entreprise sont confidentielles et incomplètes, les bases publiques décrivent les perturbations sans les décisions, et les études par simulation ne rapportent qu'un petit nombre de cas \citep{ramzy2022,do2024}. Une difficulté plus profonde s'y ajoute. Une crise réelle n'est vécue qu'une fois, sous une seule gestion, et l'on ignore ce qu'aurait donné une autre réaction. Comme les gestionnaires réagissent d'autant plus fort que la crise est grave, l'effet observé d'une décision se mêle à celui du choc.

Ce travail lève cette difficulté par le rejeu contrefactuel : chaque crise simulée est rejouée à l'identique, sans réaction puis sous plusieurs politiques de gestion, si bien que toute différence de trajectoire tient à la seule décision. Aucune base existante n'offre cette propriété. C'est elle qui permet à un modèle causal de préconiser une décision, et pas seulement de constater qu'une décision et un résultat vont souvent ensemble.

La démarche comporte trois étapes (Fig.~\ref{fig:d1}). La première, Simuler, produit des crises étiquetées, c'est-à-dire dont on connaît exactement la cause (nature, cible, date, durée et sévérité de la perturbation) et les décisions prises en réponse. Un réseau logistique multi-échelon, ici le réseau d'origine du jumeau numérique ouvert ISOMORPH \citep{zhang2026isomorph}, y est soumis à des perturbations maîtrisées auxquelles réagit un gestionnaire simulé. La deuxième, Mesurer, résume chaque trajectoire journalière du taux de service par une courbe de résilience à sept paramètres, dont on déduit six indicateurs. La troisième, Apprendre, construit sur la base ainsi obtenue un réseau bayésien causal \citep{pearl2009} qui sert au pronostic, au diagnostic et à la préconisation.

\begin{center}
\resizebox{0.88\textwidth}{!}{% Figure 1 : démarche en trois blocs (version épurée).
% L'ancienne figure détaillée reste disponible : figures/Fig1_demarche.pdf

% --- Dimensions ajustables des blocs -------------------------------------
\def\bw{5.0cm}   % largeur des blocs
\def\bh{3.0cm}   % hauteur des blocs
% -------------------------------------------------------------------------

\begin{tikzpicture}[
  font=\small,
  bloc/.style={draw=#1, line width=1pt, rounded corners=4pt, fill=#1!6,
               minimum width=\bw, minimum height=\bh},
  titre/.style={font=\small\bfseries, text=#1!80!black},
  etape/.style={draw=black!25, rounded corners=2pt, fill=white,
                text width=3.6cm, align=center, minimum height=0.62cm,
                inner sep=2pt, font=\scriptsize},
  fl/.style={-{Stealth[length=2.4mm]}, line width=1.1pt, black!55},
  pt/.style={-{Stealth[length=1.6mm]}, black!35}]

  % ---- Bloc 1 ----------------------------------------------------------
  \node[bloc=polreac] (B1) at (0,0) {};
  \node[titre=polreac, rotate=90, anchor=north]
        at ([xshift=7pt]B1.west) {1. Simuler};
  \node[etape] (a1) at (0.40, 0.85) {Réseau et perturbation étiquetée};
  \node[etape] (a2) at (0.40, 0.00) {Crise rejouée sans réaction et sous 3 profils de gestion};
  \node[etape] (a3) at (0.40,-0.85) {Évolution du taux de service $\mathrm{IP}(t)$};
  \draw[pt] (a1) -- (a2); \draw[pt] (a2) -- (a3);

  % ---- Bloc 2 ----------------------------------------------------------
  \node[bloc=palmec] (B2) at (5.6,0) {};
  \node[titre=palmec, rotate=90, anchor=north]
        at ([xshift=7pt]B2.west) {2. Mesurer};
  \node[etape] (b1) at (5.95, 0.45) {Courbes de résilience à 7 paramètres associées};
  \node[etape] (b2) at (5.95,-0.45) {6 indicateurs de résilience du réseau};
  \draw[pt] (b1) -- (b2);

  % ---- Bloc 3 ----------------------------------------------------------
  \node[bloc=polprud] (B3) at (11.2,0) {};
  \node[titre=polprud, rotate=90, anchor=north]
        at ([xshift=7pt]B3.west) {3. Apprendre};
  \node[etape] (c1) at (11.55, 0.45) {Réseau bayésien causal};
  \node[etape] (c2) at (11.55,-0.45) {Pronostic, diagnostic et préconisation};
  \draw[pt] (c1) -- (c2);

  % ---- Enchaînement ---------------------------------------------------
  \draw[fl] (B1.east) -- (B2.west);
  \draw[fl] (B2.east) -- (B3.west);
\end{tikzpicture}}
\captionof{figure}{Démarche en trois blocs : simuler des crises sous quatre politiques, mesurer leur résilience, apprendre un réseau bayésien causal.}\label{fig:d1}
\end{center}

L'article apporte un générateur de crises étiquetées et rejouées de manière contrefactuelle, applicable à tout réseau multi-échelon, qui fait de la couverture de stock un facteur d'étude plutôt qu'un obstacle. Il propose une caractérisation de la résilience dont les paramètres se lisent directement en termes logistiques. Il construit enfin un modèle causal dont la structure suit l'ordre temporel d'une crise et où la politique de gestion est une intervention, ce qui rend son effet sur la résilience identifiable.

La section 2 situe ce travail dans la littérature. La section 3 décrit la démarche, la section 4 en présente les résultats, discutés en section 5 avant la conclusion. Les développements techniques sont réunis dans un matériel supplémentaire.
